\subsection{The correspondence between subfields and Lie algebras of derivations}

We denote by E a field and by g the set of all derivations of E into itself. We recall that g is a vector space over E, with operations defined by

\[
\left(\mathrm{D} + \mathrm{D} ^ {\prime}\right) (x) = \mathrm{D} (x) + \mathrm{D} ^ {\prime} (x), \quad (a \mathrm{D}) (x) = a. \mathrm{D} (x)\tag{6}
\]

for D, D' in g and a, x in E. Moreover, if D and D' are derivations of E into E, the same is true of \([D, D'] = DD' - D'D\) (III, p.~554). Finally, Leibniz's formula (III, p.~556) gives

\[
\mathrm{D} ^ {p} (x y) = x \cdot \mathrm{D} ^ {p} (y) + \sum_ {j = 1} ^ {p - 1} \binom {p} {j} \mathrm{D} ^ {j} (x) \mathrm{D} ^ {p - j} (y) + \mathrm{D} ^ {p} (x) \cdot y
\]

(x, y in E); since the binomial coefficients \(\binom{p}{j}\) for \(1 \leqslant j \leqslant p - 1\) are divisible by p (V, p.~4, Lemma 1), we see that \(D^{p}\) is a derivation of E into E. We immediately find the relation (for \(a, a' \in E\) )

\[
\left[ a \mathrm{D}, a ^ {\prime} \mathrm{D} ^ {\prime} \right] = a a ^ {\prime} \cdot \left[ \mathrm{D}, \mathrm{D} ^ {\prime} \right] + \left(a \mathrm{D} \left(a ^ {\prime}\right)\right) \cdot \mathrm{D} ^ {\prime} - \left(a ^ {\prime} \mathrm{D} ^ {\prime} (a)\right) \cdot \mathrm{D}.\tag{7}
\]

In particular, the mapping \((\mathrm{D},\mathrm{D}^{\prime})\mapsto [\mathrm{D},\mathrm{D}^{\prime}]\) of \(\mathfrak{g}\times \mathfrak{g}\) into \(\mathfrak{g}\) is \(E^p\) -linear.

We denote by \(\mathcal{C}\) the set of subfields K of E such that \(E^p \subset K\) and \([E:K]\) is finite; for each \(K \in \mathcal{C}\) we denote by \(\mathfrak{g}(K)\) the set of K-derivations of E. Further, we denote by \(\mathcal{L}\) the set of vector subspaces \(\mathfrak{h}\) of \(\mathfrak{g}\) of finite dimension over E and such that \([D,D'] \in \mathfrak{h}\) and \(D^p \in \mathfrak{h}\) for any D, \(D'\) in \(\mathfrak{h}\) ; for each \(\mathfrak{h} \in \mathcal{L}\) we denote by \(I(\mathfrak{h})\) the set of \(x \in E\) such that \(D(x) = 0\) for all \(D \in \mathfrak{h}\) .

THEOREM 3 (Jacobson). --- The mappings \(\mathbf{K} \mapsto \mathfrak{g}(\mathbf{K})\) and \(\mathfrak{h} \mapsto \mathrm{I}(\mathfrak{h})\) are bijections of \(\mathcal{C}\) onto \(\mathcal{L}\) and of \(\mathcal{L}\) onto \(\mathcal{C}\) respectively, which are mutually inverse. If \(\mathbf{K} \in \mathcal{C}\) and \(\mathfrak{h} \in \mathcal{L}\) correspond to each other, then \([\mathrm{E}:\mathrm{K}] = p^{[\mathfrak{h}:\mathrm{E}]}\) .

The proof requires several preliminary lemmas.

Lemma 1. --- Let L be a field, V a vector space over L and u an endomorphism of V such that \(u^p = u\) . For each \(i \in \mathbf{F}_p\) let \(\mathbf{V}_i\) be the kernel of \(u - i\) . Then we have

\[
\mathbf {V} = \bigoplus_ {i \in \mathbf {F} _ {p}} \mathbf {V} _ {i}.
\]

For each \(i \in \mathbf{F}_p\) denote by \(P_i(\mathbf{X})\) the polynomial \(-\prod_{j \neq i} (\mathbf{X} - j)\) . We have

\[
(\mathbf {X} - i) \mathrm{P} _ {i} (\mathbf {X}) = \mathbf {X} - \mathbf {X} ^ {p}\tag{8}
\]

by Formula (2) (V, p.~94). On deriving the quoted formula (2), we find

\[
\sum_ {i \in \mathbf {F} _ {p}} \mathrm{P} _ {i} (\mathrm{X}) = 1.\tag{9}
\]

The formulae (8) and (9) show that the endomorphism \(\mathbf{P}_i(u)\) of \(\mathbf{V}\) maps \(\mathbf{V}\) into \(\mathbf{V}_i\) and that \(\sum_{i\in \mathbf{F}_p}\mathbf{P}_i(u) = 1\) , whence \(\mathbf{V} = \sum_{i\in \mathbf{F}_p}\mathbf{V}_i\) . It remains to show that the sum is

direct. For each \(i \in \mathbf{F}_p\) let \(v_i \in \mathbf{V}_i\) ; it is clear that \(P_i(u)\) annihilates \(v_j\) for all \(j \neq i\) and we have \(P_i(u) v_i = av_i\) with \(a = -\prod_{n \in \mathbf{F}_p^*} n \neq 0\) . The relation \(\sum_{i \in \mathbf{F}_p} v_i = 0\)

thus implies \(v_{i} = 0\) for all \(i \in \mathbf{F}_p\) and so the sum is direct.

Lemma 2. --- Let D be a derivation of E such that \(D^{p} = D\) and let K be the kernel of D. Suppose that there exists a non-zero x in E such that \(D(x) = x\) . Then K is a subfield of E containing \(E^{p}\) and we have \([E : K] = p\) .

It is clear that K is a subfield of E containing \(E^{p}\) . Denote by \(K_{i}\) the kernel of D - i for \(i \in F_{p}\) . We have \(K_{0} = K\) and Lemma 1 shows that \(E = \bigoplus_{i \in F_{p}} K_{i}\) . Let \(i \in F_{p}\) and let u be in \(K_{i}\) ; we have

\[
\mathrm{D} (x u) = \mathrm{D} (x). u + x. \mathrm{D} (u) = x u + x (i u) = (i + 1) x u
\]

whence \(xu \in K_{i+1}\) . Since x is non-zero, multiplication by x is an automorphism of the vector K-space E which maps \(K_i\) onto \(K_{i+1}\) for all \(i \in F_p\) . Since \([K_0 : K] = 1\) , we have \([K_i : K] = 1\) for all \(i \in F_p\) whence \([E : K] = p\) .

Lemma 3. --- Let \(\mathfrak{h} \in \mathcal{L}\) be of dimension \(s\) over E. Then I(h) belongs to \(\mathcal{C}\) and we have {[}E : I(h){]} = \(p^s\) .

It is clear that \(I(\mathfrak{h})\) is a subfield of E containing \(E^{p}\) . For each \(x \in E\) let \(f_{x}\) be the E-linear form \(D \mapsto D(x)\) on h. Since the intersection of the kernels of these linear forms is equal to 0, they generate the vector space dual to h (II, p.~301, Th. 7); hence there exist \(x_{1}, \ldots, x_{s}\) in E such that the linear forms \(f_{x_{1}}, \ldots, f_{x_{s}}\) form a basis of this dual. Let \((\Delta_{1}, \ldots, \Delta_{s})\) be the basis of h characterized by \(\Delta_{i}(x_{j}) = f_{x_{j}}(\Delta_{i}) = \delta_{ij}\) . Put \(D_{i} = x_{i}\Delta_{i}\) , then \((D_{1}, \ldots, D_{s})\) is a basis of h over E and we have \(D_{i}(x_{j}) = x_{i}\delta_{ij}\) . The derivations \(D_{i}^{p} - D_{i}\) and \([D_{i}, D_{j}]\) for i, j = 1, \ldots, s belong to h and annihilate \(x_{1}, \ldots, x_{s}\) ; we thus have

\[
\mathrm{D} _ {i} ^ {p} = \mathrm{D} _ {i}, \quad [ \mathrm{D} _ {i}, \mathrm{D} _ {j} ] = 0.\tag{10}
\]

For \(i\) in the range 0 to \(s\) let us denote by \(\mathbf{K}_i\) the intersections of the kernels of the derivations \(\mathrm{D}_j\) for \(1 \leqslant j \leqslant i\) . Then \(\mathbf{K}_i\) is a subfield of E and we have

\[
\mathrm{E} = {\mathrm{K}} _ {0} \supset \mathrm{K} _ {1} \supset \dots \supset \mathrm{K} _ {s - 1} \supset \mathrm{K} _ {s} = \mathrm{I} (\mathfrak {h}).
\]

Let \(i\) be between 0 and \(s-1\) ; then \(\mathbf{K}_{i}\) is stable under \(\mathbf{D}_{i+1}\) because \(\mathbf{D}_{i+1}\) commutes with \(\mathbf{D}_{1}, \ldots, \mathbf{D}_{i}\) . Moreover we have \(\mathbf{D}_{i+1}^{p} = \mathbf{D}_{i+1}, \mathbf{D}_{i+1}(x_{i+1}) = x_{i+1} \neq 0\) and \(x_{i+1} \in \mathbf{K}_{i}\) . Hence Lemma 1 implies that \([\mathbf{K}_{i}: \mathbf{K}_{i+1}] = p\) , whence finally \([\mathbf{E}: \mathbf{K}] = [\mathbf{K}_{0}: \mathbf{K}_{s}] = p^{s}\) .

We now come to the proof of the theorem. Let \(\mathfrak{h} \in \mathcal{L}\) be of dimension \(s\) over \(\mathbf{E}\) and put \(\mathbf{K} = \mathrm{I}(\mathfrak{h})\) ; then \([\mathbf{E}:\mathbf{K}] = p^s\) by Lemma 3, whence \([\Omega_{\mathbf{K}}(\mathbf{E}):\mathbf{E}] = s\) by \(\mathcal{L}\) (V, p.~103,

Th. 2, c) (V, p.~103). By the universal property of the module of differentials the mapping \(u \mapsto u \circ d_{\mathrm{E} / \mathrm{K}}\) is an isomorphism of the dual of \(\Omega_{\mathrm{K}}(\mathrm{E})\) onto \(\mathfrak{g}(\mathbf{K})\) , hence \([\mathfrak{g}(\mathbf{K}) : \mathrm{E}] = s\) . Now we have \([\mathfrak{h} : \mathrm{E}] = s\) and \(\mathfrak{h} \subset \mathfrak{g}(\mathbf{K})\) , whence \(\mathfrak{h} = \mathfrak{g}(\mathbf{K})\) , that is \(\mathfrak{h} = \mathfrak{g}(\mathrm{I}(\mathfrak{h}))\) .

Conversely, for every field \(\mathbf{K} \in \mathcal{C}\) it is clear that \(\mathfrak{g}(\mathbf{K})\) belongs to \(\mathcal{L}\) (V, p.~103, Th. 2, c)) . If \(x\) belongs to \(\mathrm{I}(\mathfrak{g}(\mathbf{K}))\) , we have \(u(d_{\mathrm{E/K}} x) = 0\) for each linear form \(u\) on \(\Omega_{\mathrm{K}}(\mathrm{E})\) , whence \(d_{\mathrm{E/K}} x = 0\) and so finally \(x \in \mathbf{K}\) by Cor. 2 of Prop. 5 (V, p.~102). Thus we have \(\mathbf{K} = \mathrm{I}(\mathfrak{g}(\mathbf{K}))\) .

Remarks. --- 1) The mutually inverse bijections \(\mathbf{K} \mapsto \mathfrak{g}(\mathbf{K})\) and \(\mathfrak{h} \mapsto \mathrm{I}(\mathfrak{h})\) are inclusion reversing; therefore \(\mathfrak{h} \mapsto \mathrm{I}(\mathfrak{h})\) is an isomorphism of the ordered set \(\mathcal{L}\) onto the ordered set opposite to \(\mathcal{C}\) . Hence we obtain the relation \(\mathrm{I}(\mathfrak{h} \cap \mathfrak{h}') = \mathrm{E}^p(\mathrm{I}(\mathfrak{h}), \mathrm{I}(\mathfrak{h}'))\) for \(\mathfrak{h}, \mathfrak{h}'\) in \(\mathcal{L}\) , because \(\mathfrak{h} \cap \mathfrak{h}'\) is the largest element of \(\mathcal{L}\) contained both in \(\mathfrak{h}\) and in \(\mathfrak{h}'\) .

\begin{enumerate}
\def\labelenumi{\arabic{enumi})}
\setcounter{enumi}{1}
\tightlist
\item
  It can be shown that every subspace of finite dimension of g which is stable under the mapping \(D \mapsto D^{p}\) is also stable under the bracket operation.
\end{enumerate}

